% 人工对照底本逐节点转写；独立TikZ重绘。

\begin{tikzpicture}[x=1cm,y=-1cm,every node/.style={anchor=west,align=left,inner sep=2pt,font=\fontsize{9.5}{12.3}\selectfont},edge/.style={draw=black!65,line width=.45pt}]
\node[text width=1.456cm] (n1) at (0.000,1.056) {量刑依据};
\node[text width=3.459cm] (n2) at (2.620,0.666) {报应刑：针对已然事实};
\node[text width=1.789cm] (n3) at (7.440,0.275) {客观危害性};
\node[text width=1.456cm] (n4) at (7.440,1.056) {主观恶性};
\node[text width=8.133cm] (n5) at (2.620,1.837) {预防刑：针对未然事实——人身危险性（再犯可能性）};
\draw[edge] (n2.east)--(7.290,0.666)--(7.290,0.275)--(n3.west);
\draw[edge] (n2.east)--(7.290,0.666)--(7.290,1.056)--(n4.west);
\draw[edge] (n1.east)--(2.470,1.056)--(2.470,0.666)--(n2.west);
\draw[edge] (n1.east)--(2.470,1.056)--(2.470,1.837)--(n5.west);
\end{tikzpicture}

